% !TEX program = xelatex
% 刻舟记 —— 随手记
% 编译:latexmk -xelatex kezhouji.tex
% 版式的意思是"本子",不是"论文":A5 的小开本、楷体、宽行距,
% 不编号、不加页码、不用脚注,注释就写在句子里。
% 字体用 TeX Live 自带的开源 Fandol,换成 macOS 系统字体就把 fandol 改成 macnew。

\documentclass[UTF8, fontset=fandol, 12pt]{ctexart}

\usepackage[a5paper, top=2.4cm, bottom=2.4cm, left=2.3cm, right=2.3cm]{geometry}
\usepackage{xcolor}
\usepackage{indentfirst}
\usepackage[hidelinks]{hyperref}

\hypersetup{pdftitle={刻舟记}, pdfsubject={随笔}}

% ——— 版式:松一点 ———
\linespread{1.6}\selectfont
\setlength{\parskip}{0.45em}
\setlength{\parindent}{2em}
\pagestyle{empty}

% 淡一号的小字,用来写题记和句子里的注
\newcommand{\dan}[1]{{\color{black!55}#1}}

% 段落之间换气,不编号,也不画记号
\newcommand{\huanqi}{\par\vspace{2.1em}}

\begin{document}

\kaishu

\noindent{\zihao{3} 刻舟记}

\vspace{0.5em}
\noindent\dan{\zihao{-6} 二〇二六年八月}

\vspace{1.6em}

\noindent\dan{\zihao{-6} 楚人有涉江者,其剑自舟中坠于水,遽契其舟曰:“是吾剑之所从坠。”舟止,从其所契者入水求之。舟已行矣,而剑不行,求剑若此,不亦惑乎?
——《吕氏春秋·察今》}

\vspace{1.8em}

小时候读到这一段,觉得世上再没有比这个楚人更笨的人了。剑掉在江心,他却在船舷上刻一道记号,等船靠了岸,再顺着那道刻痕往水里摸。老师说这叫不知变通,同学们哄堂大笑,我也跟着笑得很大声。

那时我们都以为,自己将来一定是那个划船的人,而不会是那个刻舟的人。

\huanqi

长大以后才慢慢明白,故地重游本身,就是一场刻舟求剑。

我们回到旧巷口、旧校门、旧车站,以为凭着记忆里的那道刻痕,就能在原处打捞起什么。可是水一直在流。站着的人越站越久,终于发现自己并不是在找剑,只是在原地徘徊。

年年岁岁花相似,可只有那年胜过年年。陆游写“城南小陌又逢春,只见梅花不见人”,那是沈园,他与唐琬重逢又诀别的地方,诗写在唐琬去世四十多年之后——他分明知道人已经不在了,却还是年年都去。赵嘏写“同来望月人何处?风景依稀似去年”,月色、江水、栏杆都还在原处,只有那个一同来的人不在了。

\huanqi

我们笑楚人愚钝,可自己不也一次次在记忆的船舷上刻下记号吗?

明知流水早已东去,仍忍不住回到那个路口、那家小店、那个黄昏里曾经并肩看过的晚霞。理智告诉我们,人不能两次踏进同一条河流;可情感偏要固执地相信——万一呢?万一他恰好也从这里经过呢?

这样的执拗,在旁人看来是痴,于自己却是仅剩的仪式。因为除了刻舟,我们竟找不到别的方式,来证明那段时光真实地存在过。

\huanqi

于是有人反驳得很清醒:刻舟者与那个楚人并无分别,唯有向前走,才能真正渡河。

这话是对的,却也残酷。它默认所有的河都一样宽,所有的行李都一样轻。有些人并非不懂道理,只是有些河对他们来说太沉了——沉到放下本身,就成了一种背叛。

向前走当然能到达彼岸。可彼岸没有想见的人,没有十七岁的月光。有时候我们不是过不去河,是不想过河。守着那把早已沉入江底的剑的影子,至少还能对自己说一句:我没有忘记。

\huanqi

可是江河从不因谁的停留而倒流。

我们终究要明白:剑在沉入水底的那一刻,就已经不再是一把剑了。它生锈,它被泥沙一层层裹住,一寸一寸地成为河床的一部分——成为托住整条河的东西,成为水草的养分,成为我们再往前走时,脚下那一点实实在在的着力。

刻痕会模糊,渡口会荒芜,船板会朽掉。只有水流的方向从未改变。

\huanqi

所以渡河从来不是遗忘。渡河是带着那把沉剑的重量,去往更开阔的水域。

春风再一次吹过城南小陌的时候,梅花依然开得像那年一样毫无保留。它不曾为谁少开一朵,也不曾为谁多开一朵;它只是照着自己的节令开,像水只是照着自己的方向流。

到那时我们才懂得:每一次刻舟,其实都是一次积蓄;每一次回望,都是在把力气悄悄地还给我们。那些沉进时光长河里的人和事,早已成为我们身体的一部分,在水面之下沉默地托举,把我们送到一个又一个——当年不敢想象的彼岸。

\end{document}
